% !TeX root = ../../main.tex

\label{keliyunyi}
在修井作业过程中，井筒内原有颗粒物以及作业过程中产生的新生颗粒物会显著影响井筒工况。颗粒物的存在不仅可能降低井下工具的工作效率，还可能诱发卡钻等重大安全事故。颗粒在井筒内的悬浮状态及运移特性受多种因素共同影响，主要包括颗粒粒径与密度、井眼尺寸与井斜角度，以及洗井液密度、粘度和流速等。本章基于流体动力学与颗粒动力学理论，针对砂砾、铁屑、结垢物和蜡质等典型颗粒物在井筒内的运移规律开展系统研究。通过构建Fluent-EDEM双向耦合仿真模型及多相流模型，定量分析流体参数（流速、粘度）和井筒结构参数（井斜角度）对固体颗粒运移轨迹、沉积速率及空间分布特征的影响规律。
%这里缺少一个相当于是对整个章节进行的总结

\section{基本理论和模拟条件}  

\subsection{液相控制模型}

对于密度恒定的不可压缩流体，假设液相质量为常数，液体流动遵循 N--S 方程。液固两相流模型中流体相的连续性方程和 N--S 方程由下式给出
\begin{equation}
	\frac{\partial(\rho_{\mathrm{l}}\epsilon_{\mathrm{l}})}{\partial t}
	+\nabla\cdot(\rho_{\mathrm{l}}\epsilon_{\mathrm{l}}\bm{v}_{\mathrm{l}})=0
\end{equation}
\begin{equation}
	\frac{\partial(\rho_{\mathrm{l}}\epsilon_{\mathrm{l}}\bm{v}_{\mathrm{l}})}{\partial t}
	+\nabla\cdot(\rho_{\mathrm{l}}\epsilon_{\mathrm{l}}\bm{v}_{\mathrm{l}}\bm{v}_{\mathrm{l}})
	=-\epsilon_{\mathrm{l}}\nabla P
	-\nabla\cdot(\epsilon_{\mathrm{l}}\bm{\tau}_{\mathrm{l}})
	+\epsilon_{\mathrm{l}}\rho_{\mathrm{l}}\bm{g}+\bm{F}_{p}
\end{equation}
式中，\(\rho_{\mathrm{l}}\) 为流体密度； \(\epsilon_{\mathrm{l}}\)和 \(P\) 分别为流体体积分数和压力；\(\bm{v}_{\mathrm{l}}\) 为流体速度；\(\bm{\tau}_{\mathrm{l}}\) 表示粘性应力张量；\(\bm{g}\) 为重力加速度；\(\bm{F}_{\mathrm{p}}\) 为液固相间作用力。其中，液固相间作用力及粘性应力张量公式为
\begin{equation}
	\bm{F}_{\mathrm{p}}
	=-\frac{\sum_{i=1}^{k_{\mathrm{c}}}\bm{F}_{i,\mathrm{p}}}{\Delta V}
\end{equation}
\begin{equation}
	\bm{\tau}_{\mathrm{l}}
	=\mu_{\mathrm{l}}\left[
	\nabla\bm{v}_{\mathrm{l}}
	+\left(\nabla\bm{v}_{\mathrm{l}}\right)^{\mathrm{T}}
	\right]
	-\frac{2}{3}\mu_{\mathrm{l}}
	\left(\nabla\cdot\bm{v}_{\mathrm{l}}\right)\bm{I}
\end{equation}
式中，\(\bm{F}_{i,\mathrm{p}}\) 为单个颗粒\(i\)受到的流体作用力；\(\mu_{\mathrm{l}}\) 为流体粘度，\(\bm{I}\) 为单位二阶张量。

流动需要考虑湍流的影响。考虑到高精度、应用简单、计算效率和通用性等需求，本文采用 \(k\)-\(\varepsilon\) 标准模型作为研究基础来计算钻井液的湍流粘度，其中 \(k\) 为湍流动能，\(\varepsilon\) 为耗散率。\(k\) 和 \(\varepsilon\) 的输运方程分别表示为：
\begin{equation}
	\frac{\partial}{\partial t}(\rho_{\mathrm{l}}k)
	+\frac{\partial}{\partial x_{i}}(\rho_{\mathrm{l}}kv_{\mathrm{l},i})
	=\frac{\partial}{\partial x_{j}}\left[
	\left(\mu+\frac{\mu_{\mathrm{t}}}{\sigma_{k}}\right)
	\frac{\partial k}{\partial x_{j}}
	\right]
	+G_{k}-\rho_{\mathrm{l}}\varepsilon-Y_{\mathrm{M}}
\end{equation}
\begin{equation}
	\frac{\partial}{\partial t}(\rho_{\mathrm{l}}\varepsilon)
	+\frac{\partial}{\partial x_{i}}(\rho_{\mathrm{l}}\varepsilon v_{\mathrm{l},i})
	=\frac{\partial}{\partial x_{j}}\left[
	\left(\mu+\frac{\mu_{\mathrm{t}}}{\sigma_{\varepsilon}}\right)
	\frac{\partial\varepsilon}{\partial x_{j}}
	\right]
	+C_{1\varepsilon}\frac{\varepsilon}{k}
	\left(G_{k}+C_{3\varepsilon}G_{b}\right)
	-C_{2\varepsilon}\rho_{\mathrm{l}}\frac{\varepsilon^{2}}{k}
\end{equation}
式中，\(v_{\mathrm{l}}\) 为液相速度，m/s；\(C_{1\varepsilon}\)、\(C_{2\varepsilon}\)、\(C_{3\varepsilon}\) 为常数，且 \(C_{1\varepsilon}=1.44\)，\(C_{2\varepsilon}=1.92\)。\(\sigma_{k}\) 和 \(\sigma_{\varepsilon}\) 分别为 \(k\) 和 \(\varepsilon\) 的湍流普朗特数，其取值为 \(\sigma_{k}=1.0\)，\(\sigma_{\varepsilon}=1.3\)。\(G_{k}\) 表示由于平均速度梯度产生的湍流动能；\(Y_{\mathrm{M}}\) 表示可压缩湍流脉动膨胀对整体耗散率的影响；\(\mu_{\mathrm{t}}\) 表示湍流粘度系数。\(G_{k}\)、\(Y_{\mathrm{M}}\)、\(\mu_{\mathrm{t}}\) 公式表示如下：
\begin{equation}
	Y_{\mathrm{M}}=2\rho_{\mathrm{l}}\varepsilon M^{2}
\end{equation}
\begin{equation}
	\mu_{\mathrm{t}}=\frac{\rho_{\mathrm{l}}C_{\mu}k^{2}}{\varepsilon}
\end{equation}

\subsection{EDEM理论基础} 
井筒中的颗粒在运移过程中产生颗粒-颗粒、颗粒-管壁之间的接触是不可避免的，该过程将产生作用力，本文采用EDEM软件进行模拟。EDEM软件的理论基础是离散元法（Discrete Element Method, DEM），DEM对上述两种接触对进行建模，求解接触中颗粒所受外力作用以及力矩。根据工况的不同，可区分干颗粒和湿颗粒两种类型，
当计算物料为干颗粒时，接触点不存在粘结的作用，EDEM软件默认采用Hertz-Mindlin基本接触模型，颗粒接触受力模型如图\ref{fig:particle_collision}所示。
\begin{figure}[htbp]
	\centering
	\subfloat[双颗粒碰撞模型示意图]{%
		\includegraphics[width=0.45\textwidth]{figure/kelipengzhuang.png}
		\label{fig:particle_collision}
	}
	%\vspace{0pt}
	\subfloat[Hertz--Mindlin 接触模式示意图]{%
		\includegraphics[width=0.55\textwidth]{figure/hertz-mindlin-contact-model.png}
		\label{fig:hertz_mindlin_contact_model}
	}
	\caption{颗粒--颗粒碰撞建模}
	\label{fig:particle_collision_and_contact_model}
\end{figure}

% \begin{figure}[htbp]
% 	\centering
% 	\includegraphics[width=0.6\textwidth]{figure/hertzian-soft-sphere-model.png}
% 	\caption{Hertzian 软球模型}
% 	\label{fig:hertzian_soft_sphere_model}
% \end{figure}


假设发生弹性接触的两球形颗粒的半径分别为 \(R_{1}\)、\(R_{2}\)，则颗粒间产生的法向接触力大小 \(F_{\mathrm{n}}\) 的计算式为
\begin{equation}
	F_{\mathrm{n}}=\frac{4}{3}E_{\mathrm{eq}}\sqrt{R_{\mathrm{eq}}}\,\delta^{3/2}
	\label{eq:normal_particle_elastic_force}
\end{equation}
式中，\(E_{\mathrm{eq}}\) 为等效杨氏模量，$\mathrm{GPa}$；\(\delta\) 为法向重叠量，$\mathrm{mm}$；\(R_{\mathrm{eq}}\) 为等效接触半径，$\mathrm{mm}$。\(E_{\mathrm{eq}}\) 和 \(R_{\mathrm{eq}}\) 的计算式为
\begin{equation}
	\begin{aligned}
		E_{\mathrm{eq}} & = \left(\frac{1 - \nu_1^2}{E_1} + \frac{1 - \nu_2^2}{E_2} \right)^{-1}\\
	R_{\mathrm{eq}} & =\left(\frac{1}{R_{1}}+\frac{1}{R_{2}}\right)^{-1}
	\end{aligned}
\end{equation}
式中，\(E_{1}\)、\(E_{2}\) 分别为颗粒1、2 的杨氏模量，$\mathrm{GPa}$；\( \nu_{1} \)、\( \nu_{2} \) 分别为颗粒 1、2 的泊松比。


颗粒间法向阻尼力 \(F_{\mathrm{n}}^{\mathrm{d}}\) 及切向阻尼力 \(F_{\mathrm{t}}^{\mathrm{d}}\) 的计算公式为
\begin{equation}
	\begin{aligned}
	F_{\mathrm{n}}^{\mathrm{d}} &= 2\sqrt{\frac{5}{6}}\beta\sqrt{S_{\mathrm{n}}m_{\mathrm{eq}}}\,v_{\mathrm{n}}^{\mathrm{rel}}\\
	F_{\mathrm{t}}^{\mathrm{d}} &= 2\sqrt{\frac{5}{6}}\beta\sqrt{S_{\mathrm{t}}m_{\mathrm{eq}}}\,v_{\mathrm{t}}^{\mathrm{rel}}
	\end{aligned}
	\label{eq:normal_shear_damping_force}
\end{equation}
%这是Tsuji阻尼模型
式中，\(S_{\mathrm{n}}\) 为法向刚度，\(S_{\mathrm{t}}\) 为切向刚度，$\mathrm{N/m}$；\(v_{\mathrm{n}}^{\mathrm{rel}}\) 为法向真实速度，\(v_{\mathrm{t}}^{\mathrm{rel}}\) 为切向真实速度，$\mathrm{m/s}$；\(m_{\mathrm{eq}}\) 为等效质量，\(\mathrm{kg}\)；\(\beta\) 为阻尼系数。法向和切向共用同一阻尼系数\(\beta\)，与颗粒碰撞前后的恢复系数\(e\)的关系为
\begin{equation}
	\beta = \frac{\ln e}{\sqrt{\ln^2 e + \pi ^2}}
\end{equation}
%恢复系数e定义为碰撞后法向相对分离速度与碰撞前法向相对接近速度之比：e= v_接近/v_分离

等效质量\(m_\mathrm{eq}\)的计算式为
\begin{equation}
	m_{\mathrm{eq}} = \left(\frac{1}{m_{1}}+\frac{1}{m_{2}}\right)^{-1}	
\end{equation}

\(S_{\mathrm{n}}\) 和\(S_{\mathrm{t}}\)的计算式为
\begin{equation}
	\begin{aligned}
		S_{\mathrm{n}} &= 2E_{\mathrm{eq}}\sqrt{R_{\mathrm{eq}}\delta}\\
		S_{\mathrm{t}} &= 8G_{\mathrm{eq}}\sqrt{R_{\mathrm{eq}}\delta}
	\end{aligned}
\end{equation}

其中，\(G_{\mathrm{eq}}\) 为等效剪切模量，Pa，其计算公式为
\begin{equation}
	G_{\mathrm{eq}}=\left(\frac{2-\nu_{1}}{G_{1}}+\frac{2-\nu_{2}}{G_{2}}\right)^{-1}
\end{equation}
式中，\(G_{1}\)、\(G_{2}\) 分别为颗粒 1、2 的剪切模量。

假设两颗粒产生接触前的速度分别为 \(\bm{v}_{1}\)、\(\bm{v}_{2}\)，\(\bm{n}\) 为产生接触时的法向单位矢量，则
\begin{equation}
		\bm{n}=\frac{\bm{r}_{1}-\bm{r}_{2}}{\left|\bm{r}_{1}-\bm{r}_{2}\right|}
\end{equation}
进而，\(v_{\mathrm{n}}^{\mathrm{rel}}\)和\(v_{\mathrm{t}}^{\mathrm{rel}}\)的计算式为
\begin{equation}
	\begin{aligned}
	v_{\mathrm{n}}^{\mathrm{rel}} &= (\bm{v}_{1}-\bm{v}_{2})\cdot\bm{n} \\
	v_{\mathrm{t}}^{\mathrm{rel}} &= \left|(\bm{v}_{1}-\bm{v}_{2}) - v_{\mathrm{n}}^{\mathrm{rel}}\bm{n} + (\bm{\omega}_{1} \times R_1\bm{n} + \bm{\omega}_{2} \times R_2\bm{n})\right|
	\end{aligned}
\end{equation}

颗粒间的切向力 \(F_{\mathrm{t}}\) 计算式为
\begin{equation}
	F_{\mathrm{t}}=S_{\mathrm{t}}\delta_{\mathrm{t}}
	\label{eq:shear_damping_force}
\end{equation}
%切向方向：接触面没有"几何上的切向重叠"这个概念。两个球在切向方向的相对运动，本质上是接触面发生了弹性剪切变形。这个变形量无法从当前几何位置反推，只能从接触开始的那一刻起，一步步累加相对切向位移。把切向接触想象成一个弹簧-阻尼-滑块串联系统：弹簧：代表接触面的弹性剪切刚度，弹簧的伸长量就是 δt，阻尼器：代表能量耗散，滑块：代表库仑摩擦极限
%这是Mindlin理论
式中，\(\delta_{\mathrm{t}}\)表示两个颗粒的切向重叠量，在两个颗粒建立接触时为零，此后时间步内进行累积
\begin{equation}
	\delta_{\mathrm{t}}^{n+1} = \delta_{\mathrm{t}}^{n} + \Delta t \cdot {v}_\mathrm{t}^\mathrm{rel}
\end{equation}

综合上述，根据式(\ref{eq:normal_particle_elastic_force})和式(\ref{eq:normal_shear_damping_force})，可得颗粒间的法向总接触力 ；根据式子(\ref{eq:normal_shear_damping_force})和式(\ref{eq:shear_damping_force})，可得颗粒间的切向总接触力。当然，切向摩擦力还要收到库伦摩擦定律的限制，在切向总接触里超过静摩擦力（用静摩擦系数\(\mu_s\)进行表征）时进行截断。

% 法向总接触力
\begin{equation}
    \bm{F}_{\mathrm{n}} = F_{\mathrm{n}}\bm{n} + F_{\mathrm{n}}^{\mathrm{d}}\bm{n} = \left(\frac{4}{3}E_{\mathrm{eq}}\sqrt{R_{\mathrm{eq}}}\,\delta^{3/2} + 2\sqrt{\frac{5}{6}}\beta\sqrt{S_{\mathrm{n}}m_{\mathrm{eq}}}\,v_{\mathrm{n}}^{\mathrm{rel}}\right)\bm{n}
    \label{eq:total_normal_force}
\end{equation}

% 切向总接触力（含库仑摩擦截断）
\begin{equation}
    \bm{F}_{\mathrm{t}} = 
    \begin{cases}
        -S_{\mathrm{t}}\bm{\delta}_{\mathrm{t}} - 2\sqrt{\dfrac{5}{6}}\beta\sqrt{S_{\mathrm{t}}m_{\mathrm{eq}}}\,\bm{v}_{\mathrm{t}}^{\mathrm{rel}}, & |\bm{F}_{\mathrm{t}}| \leq \mu_s F_{\mathrm{n}} \\[10pt]
        -\mu_s F_{\mathrm{n}} \dfrac{\bm{v}_{\mathrm{t}}^{\mathrm{rel}}}{v_{\mathrm{t}}^{\mathrm{rel}}}, & |\bm{F}_{\mathrm{t}}| > \mu_s F_{\mathrm{n}}
    \end{cases}
    \label{eq:total_shear_force}
\end{equation}

考虑滚动摩擦提高颗粒运动仿真准确性很重要，接触表面处产生的力矩可以用滚动摩擦系数进行表征，即
% 滚动摩擦扭矩（方向恒定扭矩模型）
\begin{equation}
    \bm{T}_{\mathrm{r}} = -\mu_{\mathrm{r}} F_{\mathrm{n}} R_{\mathrm{eq}} \frac{\bm{\omega}_{\mathrm{rel}}}{|\bm{\omega}_{\mathrm{rel}}|}
    \label{eq:rolling_torque}
\end{equation}
式中，\(\mu_{\mathrm{r}}\) 为滚动摩擦系数；\(\bm{\omega}_{\mathrm{rel}}\) 为接触点的相对角速度，\(\mathrm{rad/s}\)。以上给出了颗粒--颗粒接触碰撞过程的完整建模，对于颗粒--壁面的建模也同理，在此处不在赘述。

当处理湿颗粒且颗粒间存在液桥力等粘性作用时，接触模型通常选用 Hertz--Mindlin with JKR 模型，其法向弹性接触力的计算基于 Johnson--Kendall--Roberts（JKR）理论。
该模型在 Hertz--Mindlin（无滑移）模型的基础上引入了表面能参数，适用于模拟干粉、湿物料等强粘性系统。

JKR 模型的法向力取决于法向重叠量 $\alpha$ 和表面能 $\gamma$，其表达式为：
\begin{equation}
    F_{\mathrm{JKR}}
    = -4\sqrt{\pi\gamma E_{\mathrm{eq}}\alpha^{3}}
    + \frac{4E_{\mathrm{eq}}\alpha^{3}}{3R_{\mathrm{eq}}}
    \label{eq:JKR_force}
\end{equation}

对应的相对接近距离（即法向重叠量）$\delta$ 为：
\begin{equation}
    \delta = \frac{\alpha^{2}}{R_{\mathrm{eq}}}
    - \sqrt{\frac{4\pi\gamma\alpha}{E_{\mathrm{eq}}}}
    \label{eq:JKR_overlap}
\end{equation}

式中：$\alpha$ 为考虑粘附力时的接触面半径；$\gamma$ 为表面能（单位：J/m$^{2}$）。

当 $\gamma = 0$ 时，粘附项消失，JKR 法向力退化为经典 Hertz--Mindlin 法向力：
\begin{equation}
    F_{\mathrm{JKR}} = F_{\mathrm{Hertz}}
    = \frac{4}{3}E_{\mathrm{eq}}\sqrt{R_{\mathrm{eq}}}\,\delta^{3/2}
    \label{eq:Hertz_force}
\end{equation}

该模型即使在颗粒未发生物理接触的情况下，仍能提供吸引力（凝聚力）。颗粒间存在非零凝聚力的最大间隙 $\delta_{\mathrm{c}}$ 定义为：
\begin{equation}
    \delta_{\mathrm{c}} = \sqrt[3]{\frac{4\pi\alpha_{\mathrm{c}}}{E^{\mathrm{eq}}}}
    + \frac{\alpha_{\mathrm{c}}^{2}}{R^{\mathrm{eq}}}
    \label{eq:critical_gap}
\end{equation}

其中临界接触半径 $\alpha_{\mathrm{c}}$ 为：
\begin{equation}
    \alpha_{\mathrm{c}} = \left(\frac{9\pi\gamma R^{\mathrm{eq}2}}{2E^{\mathrm{eq}}}\right)^{\frac{1}{3}}
    \label{eq:critical_radius}
\end{equation}

当 $\delta < \delta_{\mathrm{c}}$ 时，模型返回零力。最大凝聚力（即拉脱力，pull-off force）发生在颗粒未实际接触且间隙小于 $\delta_{\mathrm{c}}$ 时，其表达式为：
\begin{equation}
    F_{\mathrm{c}} = -\frac{3}{2}\pi\gamma R^{\mathrm{eq}}
    \label{eq:pull_off_force}
\end{equation}

在实际仿真中，可根据工况调整表面能 $\gamma$ 等参数，以表征不同含水率条件下颗粒的粘性状态。表面能对颗粒流动行为影响显著，含水率越高，所需表面能参数越大。

颗粒间接触采用 EDEM 默认的 Hertz--Mindlin 模型作为基础框架，该模型适用于描述各类形状颗粒（球形、椭球形及不规则颗粒）之间的接触行为，并考虑了颗粒变形、滑移等实际因素。

在采用离散元法进行模拟计算时，通常作如下假设：
\begin{enumerate}
    \item 颗粒被视为刚性或近似刚性实体，形变仅在接触区域以重叠量体现；
    \item 颗粒间接触形式为点接触，接触发生在很小的区域内；
    \item 颗粒间可产生部分重叠，但重叠量远小于颗粒尺寸；
    \item 在每个计算步长内，外力不能同时传播到相邻颗粒，任意颗粒上的合力由其直接接触的颗粒唯一确定。
\end{enumerate}

\subsection{液固相间作用模型}

颗粒在流场中的受力主要分为以下三类：
\begin{enumerate}
	\item 与流体--颗粒相对运动无关的力，包括重力、压力梯度力、浮力等；
	\item 与流体--颗粒相对运动有关，且力的方向沿相对运动方向，包括曳力、虚拟质量力和 Basset 力等；
	\item 与流体--颗粒相对运动有关，且力的方向垂直于相对运动方向，包括 Saffman 力、Magnus 力等。
\end{enumerate}

在液--固两相流的模拟中，正确选择曳力模型有助于提高数值模拟的准确性。常见曳力模型有 Schiller--Naumann 曳力模型、Wen--Yu 曳力模型、Symmetric 曳力模型、Gidaspow 曳力模型、Syamlal--O'Brien 曳力模型等。Gidaspow 曳力模型应用较广，适用于大多数流态化计算，是 Wen--Yu 曳力模型与 Ergun 曳力模型的结合。因此，本文研究主要采用 Gidaspow 曳力模型，如下式所示：
\begin{equation}
	\beta_{\mathrm{Gidaspow}}=
	\begin{cases}
		\dfrac{3C_{\mathrm{d}}\epsilon_{\mathrm{s}}\rho_{\mathrm{l}}\left|\bm{v}_{\mathrm{l}}-\bm{v}_{\mathrm{s}}\right|}{4d_{\mathrm{s}}}\epsilon_{\mathrm{l}}^{-2.65}, & \epsilon_{\mathrm{l}}>0.8 \\[10pt]
		\dfrac{\epsilon_{\mathrm{s}}(1-\epsilon_{\mathrm{l}})\mu_{\mathrm{l}}}{\epsilon_{\mathrm{l}}d_{\mathrm{s}}^{2}}
		+1.75\dfrac{\epsilon_{\mathrm{s}}\rho_{\mathrm{l}}}{d_{\mathrm{s}}}
		\left|\bm{v}_{\mathrm{l}}-\bm{v}_{\mathrm{s}}\right|, & \epsilon_{\mathrm{l}}\leq 0.8
	\end{cases}
\end{equation}
其中，\(C_{\mathrm{d}}\) 为阻力系数，其计算公式如下：
\begin{equation}
	C_{\mathrm{d}}=
	\begin{cases}
		\dfrac{24}{Re_{\mathrm{s}}}\left(1+0.15Re_{\mathrm{s}}^{0.687}\right), & Re_{\mathrm{s}}\leq 1000 \\[8pt]
		0.44, & Re_{\mathrm{s}}>1000
	\end{cases}
\end{equation}

固相雷诺数表示为
\begin{equation}
	Re_{\mathrm{s}}=\frac{d_{\mathrm{s}}\epsilon_{\mathrm{l}}\rho_{\mathrm{l}}
		\left|\bm{v}_{\mathrm{l}}-\bm{v}_{\mathrm{s}}\right|}{\mu_{\mathrm{l}}}
\end{equation}

升力包括剪切升力（Saffman）和旋转升力（Magnus），垂直于颗粒和流体之间的相对速度方向。施加在岩屑颗粒 \(p\) 上的剪切升力（Saffman）表现形式为
\begin{equation}
	\bm{F}_{\mathrm{S}}
	=C_{\mathrm{LS}}\frac{\rho_{\mathrm{l}}\pi}{8}d_{p}^{3}
	\left|\bm{v}_{\mathrm{l}}-\bm{v}_{p}\right|
	\sqrt{\omega_{\mathrm{l}}}
\end{equation}
其中，\(\omega_{\mathrm{l}}\) 为流体速度的旋度，\(C_{\mathrm{LS}}\) 为剪切升力系数，可以表示为
\begin{equation}
	C_{\mathrm{LS}}=\frac{4.1126}{Re_{\mathrm{s}}^{0.5}}f\left(Re_{\mathrm{HB}},Re_{\mathrm{s}}\right)
\end{equation}
\begin{equation}
	f\left(Re_{\mathrm{HB}},Re_{\mathrm{s}}\right)
	=
	\begin{cases}
		\left(1-0.3314\sqrt{\beta}\right)
		e^{-0.1Re_{\mathrm{s}}}
		+0.3314\sqrt{\beta}, & Re_{\mathrm{s}}\leq 40 \\[6pt]
		0.0524\sqrt{\beta Re_{\mathrm{s}}}, & Re_{\mathrm{s}}>40
	\end{cases}
\end{equation}
其中
\begin{equation}
	\beta=\frac{0.5Re_{\mathrm{s}}}{Re_{\mathrm{HB}}},\qquad 0.005<\beta<0.4
\end{equation}

剪切流的雷诺数为
\begin{equation}
	Re_{\mathrm{s}}=\frac{\rho_{\mathrm{l}}d_{p}^{2}\omega_{\mathrm{l}}}{\mu_{\mathrm{l}}}
\end{equation}

作用在颗粒 \(p\) 上的旋转升力（Magnus）的计算公式由 Oesterle 给出：
\begin{equation}
	\bm{F}_{\mathrm{M}}
	=C_{\mathrm{LM}}\frac{\rho_{\mathrm{l}}\pi}{8}d_{p}^{3}
	\left|\bm{v}_{\mathrm{l}}-\bm{v}_{p}\right|
	\frac{\left|\bm{\Omega}\times
		\left(\bm{v}_{\mathrm{l}}-\bm{v}_{p}\right)\right|}
	{\left|\bm{\Omega}\right|}
\end{equation}
其中，\(C_{\mathrm{LM}}\) 为剪切升力系数，可以表示为：
\begin{equation}
	C_{\mathrm{LM}}
	=0.45+\left(\frac{Re_{r}}{Re_{\mathrm{HB}}}-0.454\right)
	\exp\left(-0.5684Re_{r}Re_{\mathrm{HB}}^{-0.3}\right)
\end{equation}

颗粒拟温度的物理意义表示颗粒脉动的脉动能，能够在微尺度上反映颗粒脉动速度的剧烈程度，表征颗粒运动的无序程度，是颗粒相相互碰撞作用以及颗粒相和流体相相互作用的结果。颗粒拟温度的定义为
\begin{equation}
	\theta=\frac{1}{3N}\left[
	\sum_{j=1}^{N}C_{xj}^{2}
	+\sum_{j=1}^{N}C_{yj}^{2}
	+\sum_{j=1}^{N}C_{zj}^{2}
	\right]
\end{equation}


Fluent与EDEM耦合仿真的具体实施步骤如下：首先，构建颗粒物模型并设置相关参数，基于EDEM平台按照实际物理尺寸完成颗粒注入；其次，构建三维井筒、管柱及环空的几何模型，在Fluent中建立流体仿真模型，设置流体域边界条件；随后，建立Fluent与EDEM之间的双向耦合接口，模拟真实井筒工况下流速、粘度等关键参数对颗粒物运移轨迹、沉积速率及空间分布特征的影响，具体操作流程如图\ref{fig:flowchart}所示；最后，利用后处理模块对仿真结果进行定量分析，包括颗粒分布统计和动力学参数提取。

    \begin{figure}[htbp]
        \centering
        \begin{tikzpicture}[node distance=2cm] % 节点间距为2cm
            % 1. 流场初始化（起始节点）
            \node (init) [startstop] {流场初始化};
            % 2. 连续相流场计算（处理节点）
            \node (flow) [process, below of=init] {连续相流场计算};
            % 3. 颗粒模型注入（处理节点）
            \node (inject) [process, below of=flow] {颗粒模型注入};
            % 4. 颗粒时间步长计算（处理节点）
            \node (time1) [process, below of=inject] {颗粒时间步长计算};
            % 6. 颗粒轨迹计算（处理节点）
            \node (trajectory) [process, below of=time1] {颗粒轨迹计算};
            % 7. 更新连续相源项（处理节点）
            \node (update) [process, below of=trajectory] {更新连续相源项};
            % 8. 判断是否收敛（决策节点）
            \node (converge) [decision, below of=update] {是否收敛};
            % 9. 判断是否计算完成（决策节点）
            \node (complete) [decision, below of=converge] {是否计算完成};
            % 10. 计算结果（终止节点）
            \node (result) [startstop, below of=complete] {计算结果};
            % 11. 更新下一步时间数据（右侧辅助节点）
            \node (nextstep) [process, right of=complete, xshift=4cm] {更新下一步时间数据};

            % 绘制箭头连线
            \draw [arrow] (init) -- (flow);                          % 流场初始化 → 连续相流场计算
            \draw [arrow] (flow) -- (inject);                        % 连续相流场计算 → 颗粒模型注入
            \draw [arrow] (inject) -- (time1);                       % 颗粒模型注入 → 颗粒时间步长计算
            \draw [arrow] (time1) -- (trajectory);                   % 颗粒时间步长计算 → 颗粒轨迹计算
            \draw [arrow] (trajectory) -- (update);                  % 颗粒轨迹计算 → 更新连续相源项
            \draw [arrow] (update) -- (converge);                    % 更新连续相源项 → 判断是否收敛

            % 收敛判断分支（未收敛→返回连续相流场计算）
            \draw [arrow] (converge.west) -- ++(-1,0) |- (flow.west) node[midway, left] {否};

            % 收敛后→计算完成判断
            \draw [arrow] (converge) -- (complete) node [midway, right] {是};

            % 计算完成判断分支（未完成→更新下一步时间数据→返回连续相流场计算）
            \draw [arrow] (complete.east) -- (nextstep.west) node[midway, above] {否};
            \draw [arrow] (nextstep.north) |- (flow.east);

            % 计算完成→输出结果
            \draw [arrow] (complete) -- (result) node[midway, right] {是};
        \end{tikzpicture}
        \caption{颗粒运移模型计算流程图}
        \label{fig:flowchart}  
    \end{figure}


在井筒内颗粒物运移特性的研究中，基于对颗粒动力学行为控制作用的大小，选择流体流速 \(v_{\mathrm{c}}\)、流体粘度 \(\mu\) 和井斜角 \(\theta\) 作为主要控制变量。首先，流体流速 \(v_{\mathrm{c}}\) 是决定颗粒悬浮与输运能力的关键参数，流速直接影响流体的动能和湍流强度，进而调控颗粒的启动阈值、悬浮状态及沉积风险。其次，流体粘度 \(\mu\) 通过改变流体的粘滞阻力，显著影响颗粒的沉降速度及与流体的相互作用，高粘度流体可增强颗粒携带能力，但可能增加流动阻力。最后，井斜角 \(\theta\) 作为几何特征的核心变量，通过改变重力在井筒轴线方向的分量，直接影响颗粒的受力平衡：在斜井或水平井中，重力分量的减弱会加剧颗粒沉降的方位选择性，导致非对称沉积。这三个因素共同构成了颗粒运移的“动力--阻力--几何”三元耦合机制，能够系统表征复杂井筒环境下颗粒迁移的物理本质，为优化携粒效率、预测床层形成提供理论支撑。考虑到颗粒属性的不同会对其运移行为产生显著影响，于是，下文分别研究砂砾、铁屑、结垢物和蜡颗粒的运移和分布规律。

